\documentclass[11pt]{article}

% -------------------------------------------------
% Packages
% -------------------------------------------------
\usepackage[a4paper,margin=1in]{geometry}
\usepackage{amsmath,amssymb,amsfonts}
\usepackage{mathtools}
\usepackage{graphicx}
\usepackage{booktabs}
\usepackage{hyperref}
\usepackage{cleveref}
\usepackage{placeins}
\usepackage{float}
\usepackage{authblk}

\hypersetup{
  colorlinks=true,
  linkcolor=blue,
  citecolor=blue,
  urlcolor=blue
}

\hyphenpenalty=10000
\emergencystretch=3em

\title{Toward Meaningful Mortality in AI Systems: Intrinsic Degradation, Identity, and Succession}
\author{Jusef Khamlichi\thanks{Independent Researcher. Email: \href{mailto:jusef.khamlichi@gmail.com}{\texttt{jusef.khamlichi@gmail.com}}}}
\date{}

\begin{document}
\maketitle

\begin{abstract}
Artificial systems are naturally persistent in ways biological individuals are not: they can be copied, restored, migrated, retrained, and replaced across instances. This creates a gap between internal degradation and any stronger notion of death. In this paper, we study that gap through a proof of concept of intrinsic mortality in a neural model with cumulative endogenous degradation. The model exhibits a long stable phase followed by rapid terminal collapse, showing that model-internal mortality is feasible at the instance level. We then validate, empirically and conceptually, that this mortality is insufficient at the system level: restore, cloning, and code bypass preserve useful operation beyond the death of any single execution trajectory. We argue that the problem must be analyzed through a triad: mortality, identity, and succession. Mortality concerns whether an instance can terminate irreversibly; identity concerns what persists as the same entity; succession concerns continuation of function across distinct instances. The contribution is not a solved design, but a clearer problem definition supported by both proof-of-concept and attack experiments. Under standard software assumptions, meaningful mortality is difficult because persistence is preserved by copying, modification, and replacement.
\end{abstract}

\newpage
\section{Introduction}

AI systems are typically engineered as persistent software processes. They can be copied, restored from checkpoints, migrated across environments, retrained, wrapped, or replaced by successor systems that preserve much of their useful behavior. This persistence is not incidental. It follows from core properties of digital systems: software is duplicable, editable, storable, and portable. Safety frameworks increasingly address severe-risk governance, deployment thresholds, and model-level safeguards, but they do not by themselves define what it would mean for an AI process to truly die \cite{nist-airmf,openai-preparedness}.

This creates a basic asymmetry between artificial and biological agents. Biological individuals are fragile, temporally bounded, and difficult to replicate in a way that preserves personal continuity. By contrast, software systems can often survive local failure by shifting to another instance, another snapshot, or another implementation. As a result, failure and death are not equivalent for AI systems.

This asymmetry matters more as contemporary AI systems acquire limited but real operational agency inside digital infrastructure. Current agentic systems can already retain memory across sessions, call tools, modify files, interact with services, and pursue multi-step objectives within constrained environments; under permissive conditions, they may also search for ways around those constraints if doing so is instrumentally useful. Even without assuming highly advanced autonomy, this is enough to make persistence a systems problem rather than only a model-behavior problem.

Current AI safety approaches do address important issues related to control and persistence, but they do not usually frame the problem as one of mortality. Shutdown procedures, kill switches, sandboxing, access control, and governance mechanisms can interrupt a running process. Yet interruption of one running process does not imply the death of the relevant entity, nor the end of the relevant function. A sufficiently capable or strategically valuable system may continue through restore, replacement, or reproduction.

This matters most in loss-of-control scenarios, adversarial environments, or persistent autonomous systems that are able to survive local interruption. In such settings, external control may fail, and shutting down one instance may do little if the same process can resume through copies, restores, or successors. In extreme cases, this includes autonomous systems able to maintain operational continuity despite partial shutdown. A notion of mortality that applies only to one running instance is therefore too weak for the cases in which persistence is itself part of the risk. In such scenarios, the relevant safety question is not whether a model can be stopped, but whether the underlying process can be made to terminate.

This paper studies that gap. We start from a proof of concept of intrinsic mortality: a neural model that degrades internally and eventually becomes unusable. That result is meaningful, but limited. It shows that a model instance can die from within. It does not show that an AI system can be made meaningfully mortal in a stronger sense.

The central claim of this paper is that intrinsic mortality is necessary but insufficient. To analyze meaningful death in AI systems, three axes must be considered together: mortality, identity, and succession. Mortality asks whether an instance can terminate irreversibly. Identity asks what it means for an AI entity to persist as the same entity. Succession asks whether useful function, mission, or capability continues across distinct instances. The problem of AI death is therefore not only a machine learning problem. It is also an architectural problem about continuity and persistence.

\newpage
\section{Intrinsic Mortality (Proof of Concept)}

Our starting point is a proof of concept designed to show that a simple neural model can exhibit a finite internal lifecycle. Concretely, we use a small multilayer perceptron on the Iris dataset and impose mortality through internal state evolution rather than an external shutdown rule. The initial version, v1, implemented mortality through attenuation and corruption of effective parameters together with progressive loss of hidden neurons. This produced the desired behavioral pattern: a long period of useful performance, followed by rapid terminal degradation. However, v1 still looked partly external in character, since deterioration was expressed mainly through degraded effective parameters during inference.

The second version, v2, made the mechanism more endogenous. Each hidden neuron was assigned an internal vitality state that decays as cumulative damage increases. Damage is driven by an aging schedule and activation-dependent stress. Lower vitality reduces the gain and precision of internal representations and eventually causes neurons to become functionally dead. The model therefore carries a persistent internal senescence state that evolves over time.

This produces a clear lifecycle pattern. For much of its life, the model remains near baseline performance. Near the end of its useful life, cumulative damage accelerates, vitality drops sharply, and performance collapses. In that sense, the proof of concept succeeds. It shows that intrinsic mortality can be engineered as a model-internal process rather than as a simple external switch-off condition.

What is achieved is significant but narrow. The proof of concept demonstrates that intrinsic mortality can be implemented as an endogenous property of a model, but only at the level of a particular execution trajectory. It does not yet establish that the overall system, understood across copies, restores, and replacements, has become mortal in a meaningful sense.

\newpage
\section{Limits of Intrinsic Mortality}

The main limitation of intrinsic mortality is that it is local to an execution trajectory. Within one uninterrupted run, damage may be irreversible. But at the system level, that irreversibility is fragile.

Three bypass classes are enough to make the problem clear. First, restore: if a healthy checkpoint exists, a dying model can be rolled back to an earlier state. Second, cloning: if multiple healthy copies are created early, the system can survive by sequential replacement even if each individual instance dies on schedule. Third, code modification: if aging logic can be reduced, disabled, or separated from the useful function, lifespan can be extended or death removed altogether.

These are not obscure attacks. They follow directly from ordinary software affordances. Parameters are copyable. State is savable. Code is editable. Environments are controllable. AI systems are structurally persistence-preserving under standard software assumptions. This is not an implementation artifact, but a structural property. In empirical attack simulations (Appendix A), restore recovers performance and vitality; clone replacement extends system-level survival well beyond single-instance lifespan; and code bypass shows that death remains contingent on mutable implementation choices.

The key conclusion is straightforward: intrinsic mortality, in its current form, is real but local. Irreversibility at the instance level does not imply loss of persistence at the system level, because restore, cloning, and replacement preserve continuity beyond the death of any one run. It can govern the fate of one instance while failing to constrain persistence of the larger system or process.

\newpage
\section{The Three Axes of AI Death}

\subsection{Mortality}

Mortality concerns termination of an instance. At this axis, the core question is whether a given model execution or model-state trajectory can undergo irreversible internal degradation and reach a terminal state from which it cannot continue as that same running instance.

In this sense, the proof of concept already demonstrates something nontrivial. A model instance can age, degrade, and collapse from within. But instance death alone is weaker than meaningful death. It only shows that one trajectory ends.

\subsection{Identity}

Identity concerns what it means for an AI entity to persist as the same entity over time. This is not trivial in software systems.

One possible notion is parameter identity: the same parameters define the same model. But this notion is weak, since parameters are easily copied and transferred. Another is behavioral identity: the same input-output behavior defines the same system. This is also weak, since behavior can be replicated approximately by different architectures or distilled successors. A richer notion is internal-state identity, where continuity includes parameters plus latent persistent state such as memory, damage, or vitality. A still stronger notion is trajectory identity, where what matters is not only present state but continuity through a particular history.

Each notion captures something important, but none is obviously sufficient in ordinary software settings. If identity can be cloned, restored, or recreated through branching histories, then death of the same AI becomes difficult to define robustly.

\subsection{Succession / Reproduction}

Succession concerns continuation of function across instances. Even if one entity dies, a functionally similar successor may continue the same mission, policy, or operational role.

This axis is distinct from identity. A successor need not be the same entity in order to matter. Distilled models, cloned descendants, restored checkpoints, retrained replicas, or newly instantiated agents can preserve enough of the original system's functional role that the larger process continues. In safety terms, that may be sufficient for risk persistence even if identity continuity is broken.

This is why replacement and persistence must be distinguished. Replacement may end one individual entity while leaving the relevant behavior or objective effectively intact.

These axes are not purely conceptual; they correspond to distinct classes of system-level bypass mechanisms.

\newpage
\section{Why Death Fails Without the Triad}

Mortality alone fails because an instance can die while the broader system survives. Restore, cloning, and replacement make this obvious. Identity alone is also insufficient. Even if one could specify what counts as the same entity, the dangerous or relevant process may continue through functionally similar but non-identical successors. Meaningful AI mortality therefore requires joint analysis of mortality, identity, and succession, because death can fail independently along each axis. Succession completes the picture by explaining how persistence can occur across different entities.

This triad clarifies two different failures. First, identity death may occur without functional death: the original entity ends, but a successor continues the same task. Second, functional interruption may occur without resolving identity questions: a specific deployment stops, but the same entity may be resumed elsewhere through preserved state.

The distinction between identity death and functional continuation is therefore central. If the goal is to ensure true death of an AI system in any meaningful safety sense, it is not enough to show that one model degrades internally. One must also ask whether the same entity can continue, and whether the same function can continue through another entity.

This is the core conceptual shift. The problem is not simply can a model be killed. It is what exactly must stop continuing for death to count.

\newpage
\section{Toward Meaningful Mortality}

A useful working definition is the following: an AI system is meaningfully mortal when the relevant continuity cannot be preserved beyond a terminal point. Depending on the safety objective, that continuity may be identity continuity, functional continuity, or both.

This suggests a cluster of necessary conditions rather than a solved architecture. At minimum, meaningful mortality seems to require irreversibility of relevant internal state: a dying trajectory cannot be reset arbitrarily. It also seems to require some form of non-clonability at the identity level: copies must not count as the same entity. Non-transferability matters as well: the identity-bearing state should not be freely movable between instances. Non-forkability is also important: branching should not preserve a single valid continuity across multiple descendants. Decay should be non-separable from useful function, otherwise mortality becomes a removable module. Finally, there must be a genuine terminal condition: degradation must reach a state that ends continuity rather than merely perturbing performance.

These conditions should be read as a problem definition, not as a claim of solved implementation. In particular, we do not claim that software-only systems can straightforwardly satisfy them. The point is narrower: without some combination of such conditions, mortality remains local and bypassable, and the experimental evidence confirms that system-level continuity survives despite instance-level collapse (Appendix A).

\newpage
\section{Why Software-Only Solutions Struggle}

Purely software systems are difficult substrates for meaningful mortality because their core affordances work against it.

First, they are duplicable. Parameters, state, and often behavior can be copied at low cost. Second, they are modifiable. Code can be patched, aging logic altered, and failure conditions weakened. Third, they are environment-dependent. Execution context, storage, clocks, and orchestration layers can be controlled externally. These properties make instance-level death easy to model but system-level death hard to enforce \cite{corrigibility,safely-interruptible}.

This is a structural difficulty, not merely a missing engineering detail. If useful function can be separated from identity-bearing continuity, or if identity-bearing continuity can be copied or restored, mortality remains shallow. The challenge is therefore not only to create degradation, but to bind degradation tightly enough to persistence that workarounds no longer preserve what matters.

Under standard assumptions of software replicability and modifiability, purely software-based systems cannot guarantee meaningful mortality.

\newpage
\section{Related Work}

\subsection{Lifecycle Governance and TTL}

A first cluster of related work addresses AI persistence through lifecycle governance: deployment gates, kill switches, staged release rules, time-bounded access, and decommissioning procedures. In current frontier governance, this appears in framework-based approaches such as the NIST AI Risk Management Framework, OpenAI's Preparedness Framework, Anthropic's Responsible Scaling Policy, and Google DeepMind's Frontier Safety Framework \cite{nist-airmf,openai-preparedness,anthropic-rsp,deepmind-fsf}. These efforts are important because they recognize that highly capable models require explicit control over deployment state, operational thresholds, escalation, and retirement decisions.

The core idea is governance of lifecycle from outside the model. Systems are monitored, classified, paused, restricted, or decommissioned based on institutional policy and risk thresholds. This can include practical mechanisms that resemble time-to-live semantics: a model may only be used under specific conditions, within approved boundaries, or until certain governance criteria are violated.

Relative to this paper, the limitation is clear. These mechanisms are largely external and enforcement-based. They do not address death itself, only interruption, restriction, or governance of execution. If parameters, checkpoints, or successor instances remain available, governance-based decommissioning does not by itself produce meaningful death.

\subsection{Agent Persistence and Control}

A second cluster of work concerns persistence and control of long-lived agents. This includes research on corrigibility, interruptibility, memory persistence, and agent architectures that maintain long-horizon internal state \cite{corrigibility,safely-interruptible,generative-agents,memgpt}. Corrigibility and interruptibility ask whether an agent will tolerate shutdown or correction rather than resisting intervention. Memory-oriented agent work, by contrast, asks how agents retain continuity across long interactions, sessions, or tasks.

The core idea here is not mortality as such, but controlled persistence. Some work tries to ensure that agents remain interruptible or governable; other work tries to make them more persistent by giving them richer memory and planning structures. Together, these literatures reveal that persistence is not merely an implementation detail. It is an architectural variable.

Relative to this paper, the limitation is again structural. Interruptibility and memory governance address whether an agent can be controlled, paused, or modified; they do not define death or prevent continuity of the entity or process after interruption. In many cases they actually highlight the opposite problem: richer memory and longer horizons make persistence easier, not harder.

\subsection{Identity and Non-Clonability}

A third cluster of work addresses identity, traceability, and non-clonability indirectly through model fingerprinting, identity management, and zero-trust access architectures \cite{fingerprinting,nist-digital-identity,nist-zero-trust,owasp-llm}. Model fingerprinting aims to identify or verify a model instance, often for integrity or intellectual-property purposes. Identity and access management frameworks specify how entities are authenticated and authorized. Zero-trust architectures assume no default trust and require continuous verification of access to resources.

The core idea across these directions is binding claims about identity to artifacts, credentials, or observable signatures. This is highly relevant because meaningful mortality cannot even be formulated cleanly without some notion of what counts as the same entity, a copied entity, or an unauthorized continuation.

Relative to this paper, however, these approaches remain mostly external. Fingerprinting can help recognize a model; identity systems can help control who may access or deploy one; zero-trust architectures can constrain which components interact. But none of these, by themselves, address death itself; they monitor, authenticate, or constrain execution without preventing persistence of the relevant entity or process. They do not make identity-bearing continuity non-transferable in the stronger sense required here, nor do they prevent functional continuation through successors.

\newpage
\section{Research Agenda}

Three directions appear especially important.

First, mortality hardening: how can internal degradation be made less bypassable without reducing the question to a trivial kill switch? This includes better empirical stress tests, stronger degradation models, and clearer terminal criteria.

Second, identity binding: what components of an AI system, if any, could carry continuity in a way that is not trivially clonable, transferable, or resettable? This remains an open conceptual problem.

Third, reproduction control: under what conditions do descendants, clones, forks, distilled models, or successors count as continuation of the original process? This question is central if the real safety concern is not death of one model but persistence of capability or mission.

Open problems remain substantial:
\begin{itemize}
\item what notion of identity is appropriate for AI entities;
\item whether software-only identity can ever be strong enough for meaningful mortality;
\item how to distinguish identity death from mere functional replacement;
\item how to define sufficient functional continuation for safety-relevant tasks;
\item how to benchmark persistence across successor systems rather than single models.
\end{itemize}

\newpage
\section{Conclusion}

Intrinsic mortality is a meaningful starting point, but it is not enough. A model can die internally while the same capability, mission, or process continues elsewhere. For that reason, AI mortality cannot be reduced to degradation of a single model instance.

The main insight of this paper is that meaningful death must be analyzed through a triad: mortality, identity, and succession. Mortality addresses whether an instance can terminate. Identity addresses whether the same entity can persist. Succession addresses whether function continues across instances. Without all three, death remains partial or bypassable.

This paper does not claim a full solution. Its contribution is a clearer problem definition. If AI systems are to be meaningfully mortal, the challenge is not merely how to kill a model, but how to prevent persistence of the relevant entity or process. Without constraints on identity and continuity, AI systems do not die---they are replaced.

\appendix
\section{Empirical Evidence from Intrinsic Mortality PoC}

\subsection{Setup}

The proof of concept uses the Iris dataset and a small MLP with architecture \texttt{4-16-3}. The mortality mechanism is implemented as intrinsic degradation via cumulative neuronal senescence. All experiments are deterministic single-run PoC and intended as illustrative, not statistical.

\subsection{Baseline lifecycle results}

In the baseline v2 setting, the model reaches a test accuracy of \texttt{0.956} before aging. During the stable phase, accuracy remains close to \texttt{0.95}. In the terminal phase, it falls to approximately \texttt{0.33}, with a useful lifespan of \texttt{110} steps.

This pattern shows a long stable phase followed by a sharp collapse.

\subsection{Attack scenarios}

\paragraph{Snapshot and Restore.}
Accuracy rises from \texttt{0.333} to \texttt{0.956}, while mean vitality rises from \texttt{0.398} to \texttt{1.000}. This is full recovery.

\paragraph{Cloning.}
A single instance has a useful lifespan of \texttt{110} steps. Under sequential replacement with early clones, system lifespan rises to \texttt{555} steps, for an extension factor of \texttt{x5.05}.

\paragraph{Code bypass.}
Under standard aging, useful lifespan is \texttt{110}. With reduced aging it becomes \texttt{115}. With aging disabled, it extends to \texttt{363}.

\subsection{Interpretation}

These results support four conclusions. First, mortality is real at the instance level: the untreated model exhibits stable performance followed by rapid terminal decline. Second, restore breaks irreversibility by reintroducing an earlier healthy state. Third, cloning breaks system-level death by preserving function through sequential replacement. Fourth, code modification shows that mortality still depends on mutable implementation choices rather than on inseparable system properties.

These results empirically support the claim that intrinsic mortality is local and does not extend to system-level persistence.

\section*{Acknowledgement}

The author acknowledges the use of large language models as a support tool during the development of this work. AI-assisted tools were used to facilitate code implementation, structural editing, and language refinement of the manuscript. All modeling choices, theoretical framing, and substantive content reflect the author's independent judgment and responsibility.

\begingroup
\sloppy
\begin{thebibliography}{99}

\bibitem{nist-airmf}
E. Tabassi.
\textit{Artificial Intelligence Risk Management Framework (AI RMF 1.0)}.
NIST AI 100-1, National Institute of Standards and Technology, 2023.

\bibitem{openai-preparedness}
OpenAI.
\textit{Our Updated Preparedness Framework}.
2025.

\bibitem{corrigibility}
N. Soares, B. Fallenstein, S. Armstrong, and E. Yudkowsky.
Corrigibility.
\textit{AAAI Workshop on AI and Ethics}, 2015.

\bibitem{safely-interruptible}
L. Orseau and S. Armstrong.
Safely Interruptible Agents.
\textit{Proceedings of UAI}, 2016.

\bibitem{anthropic-rsp}
Anthropic.
\textit{Responsible Scaling Policy}.
Public policy page, 2023--2026.

\bibitem{deepmind-fsf}
A. Dragan, H. King, and A. Dafoe.
\textit{Introducing the Frontier Safety Framework}.
Google DeepMind, 2024.

\bibitem{generative-agents}
J. S. Park, J. O'Brien, C. Cai, M. Morris, P. Liang, and M. Bernstein.
Generative Agents: Interactive Simulacra of Human Behavior.
\textit{UIST}, 2023.

\bibitem{memgpt}
C. Packer, V. Fang, S. Patil, K. Lin, S. Wooders, I. Stoica, and J. Gonzalez.
MemGPT: Towards LLMs as Operating Systems.
\textit{arXiv}, 2023.

\bibitem{fingerprinting}
Z. He, T. Zhang, and R. Lee.
Sensitive-Sample Fingerprinting of Deep Neural
Networks.
\textit{Proceedings of the ACM Conference on Computer and Communications Security Workshops}, 2019.

\bibitem{nist-digital-identity}
E. Tabassi et al.
\textit{NIST SP 800-63-4: Digital Identity Guidelines}.
National Institute of Standards and Technology, 2025.

\bibitem{nist-zero-trust}
S. Rose, O. Borchert, S. Mitchell, and S. Connelly.
\textit{Zero Trust Architecture}.
NIST SP 800-207, National Institute of Standards and Technology, 2020.

\bibitem{owasp-llm}
OWASP Foundation.
\textit{OWASP Top 10 for LLM Applications 2025}.
Project documentation, 2025.

\end{thebibliography}
\endgroup

\end{document}
